\subsection{Statistical guarantees}
\label{subsec:elbo_tight2}

This subsection separates three questions that are sometimes conflated.  The
first is an inference question: for fixed generative parameters, how far is
the encoder distribution from the exact posterior?  The second is an
estimation question: does maximizing an approximate likelihood recover the
generative parameter?  The third is an algorithmic question: does the
numerical procedure find a sufficiently accurate maximizer?  Tightness of the
ELBO answers the first question.  Consistency and asymptotic normality require
additional identification, empirical-process, and optimization conditions.

To include random modality subsampling without confusing it with sampling of
units, let the full observed record be
\[
 U_i:=(O_i,w_{i,O_i},y_{i,1:K},x_i).
\]
For $S\subseteq O_i$, define
\[
 u_{i,S}:=(w_{i,S},y_{i,1:K}),\qquad
 V_i(S):=(S,u_{i,S},x_i),
\]
and write
\[
 \ell_\vartheta(V_i(S)):=\log p_\vartheta(u_{i,S}\mid x_i,S),
 \qquad
 \pi_\vartheta(z\mid V_i(S))
 :=p_\vartheta(z\mid u_{i,S},x_i,S).
\]
Here $\vartheta=(\theta,\psi)$.  The exact expected-subset per-unit version
of the training objective is the composite contribution
\begin{equation}
 \bar\ell_\vartheta(U_i)
 :=\E_{S\sim r(\cdot\mid O_i)}[\ell_\vartheta(V_i(S))].
 \label{eq:composite_subset_criterion}
\end{equation}
All statements apply to a fixed subset by taking $S$ degenerate.  Nothing in
the theory distinguishes $w$ from $y$: their separate notation is
expositional, and both are components of $u_{i,S}$.  As written, the results
require ignorability of $O_i$, or directly require correct specification of
the selected subset marginals.  With informative missingness one must replace
$\ell_\vartheta$ by a correctly specified selection or observed-data
likelihood and recheck the identification and regularity conditions.
Tightness for the posterior conditional on
$u_{i,S}$ also requires the encoder to receive $u_{i,S}$ (or a sufficient
statistic).  An encoder intentionally denied $y_i$ cannot generally equal
$p_\vartheta(z_i\mid w_{i,S},y_i,x_i)$.

\subsubsection{The inference gap}

The following identity is the starting point for all of the results.

\begin{proposition}[ELBO identity and gap decomposition]
\label{prop:elbo_identity_decomposition}
For every $\vartheta$, every realization $V$, and every density $q(\cdot\mid
V)$ for which the displayed quantities are finite,
\begin{equation}
 \ell_\vartheta(V)-\mathcal L(\vartheta,q;V)
 =\mathrm{KL}\!\left(q(\cdot\mid V)\,\middle\|\,
 \pi_\vartheta(\cdot\mid V)\right)\geq 0.
 \label{eq:elbo_gap_identity_profile}
\end{equation}
Let $\mathcal Q$ be a per-observation variational family and define
\[
 \Delta_{\mathcal Q}(\vartheta,V)
 :=\inf_{q\in\mathcal Q}\mathrm{KL}
 \!\left(q\,\middle\|\,\pi_\vartheta(\cdot\mid V)\right).
\]
For an amortized encoder $q_\phi(\cdot\mid V)\in\mathcal Q$, the total gap
has the exact decomposition
\begin{align}
 \ell_\vartheta(V)-\mathcal L(\vartheta,q_\phi;V)
 &=\underbrace{\Delta_{\mathcal Q}(\vartheta,V)}_{
 \mathrm{variational\ gap}}
 \nonumber\\
 &\quad+
 \underbrace{\left\{
 \mathrm{KL}(q_\phi\|\pi_\vartheta)-
 \Delta_{\mathcal Q}(\vartheta,V)\right\}}_{
 \mathrm{amortization\ gap}}.
 \label{eq:gap_decomposition}
\end{align}
Both terms are nonnegative.  The ELBO is tight at $V$ if and only if
$q_\phi(\cdot\mid V)=\pi_\vartheta(\cdot\mid V)$ almost everywhere.
\end{proposition}

This decomposition follows \citet{cremer2018inference}.  It is pointwise:
small average reconstruction error, or a flexible encoder in isolation, does
not imply a small inference gap.

\subsubsection{One-sided comparison with the exact criterion}
\label{subsubsec:one_sided_comparison}

The most direct consistency argument uses only the lower-bound identity and
one good encoder comparator.  It does not require a uniform law over the
neural-network class.  Put
\begin{align}
 \ell_n(\vartheta)&:=P_n\bar\ell_\vartheta,\\
 L_n(\vartheta,q)&:=P_n\bar m(\vartheta,q),\\
 D_n(\vartheta,q)&:=\ell_n(\vartheta)-L_n(\vartheta,q)\\
 &=P_n\E_{S\mid O}\mathrm{KL}\{q(\cdot\mid V(S))\|\pi_\vartheta(\cdot\mid V(S))\}\geq0.
 \label{eq:normalized_elbo_comparison}
\end{align}
All optimization and gap rates below are on this per-unit scale.

\begin{theorem}[One-comparator consistency and aggressive comparison]
\label{thm:one_comparator}
Suppose that, for every fixed $\epsilon>0$, with probability tending to one,
\[
 \sup_{\operatorname{dist}(\vartheta,\vartheta_0)\geq\epsilon}
 \ell_n(\vartheta)
 \leq \ell_n(\vartheta_0)-c_\epsilon
\]
for some $c_\epsilon>0$.  If one deterministic $q_n^0\in\mathcal Q_n$
satisfies $D_n(\vartheta_0,q_n^0)=o_p(1)$ and the computed pair obeys
\[
 L_n(\widehat\vartheta_n,\widehat q_n)
 \geq\sup_{\vartheta,q\in\mathcal Q_n}L_n(\vartheta,q)-o_p(1),
\]
then $\operatorname{dist}(\widehat\vartheta_n,\vartheta_0)\to_p0$.  If also
$\ell_n(\widehat\vartheta_n)-\ell_n(\vartheta_0)=o_p(1)$, then
$D_n(\widehat\vartheta_n,\widehat q_n)=o_p(1)$.

For the sharper comparison, let $\widetilde\vartheta_n$ maximize $\ell_n$.
Suppose both estimators lie in a common identified neighborhood on which
\begin{equation}
 \ell_n(\widetilde\vartheta_n)-\ell_n(\vartheta)
 \geq c\|\vartheta-\widetilde\vartheta_n\|^2,
 \label{eq:exact_local_quadratic_separation}
\end{equation}
one (possibly data-dependent) $q_n^\dagger\in\mathcal Q_n$ satisfies
$D_n(\widetilde\vartheta_n,q_n^\dagger)\leq\rho_n$, and the global per-unit
ELBO optimization loss is at most $\epsilon_n$.  Then
\begin{align}
 \|\widehat\vartheta_n-\widetilde\vartheta_n\|
 &\leq c^{-1/2}(\rho_n+\epsilon_n)^{1/2},
 \label{eq:aggressive_parameter_comparison}\\
 D_n(\widehat\vartheta_n,\widehat q_n)
 &\leq\rho_n+\epsilon_n.
 \label{eq:aggressive_fitted_gap}
\end{align}
Consequently, $\rho_n+\epsilon_n=o_p(n^{-1})$ makes the raw ELBO estimator
root-$n$ equivalent to the exact extremum estimator.  The weaker order
$o_p(n^{-1/2})$ makes it an $o_p(n^{-1/4})$ preliminary estimator when the
exact estimator is root-$n$ regular.
\end{theorem}

The extra fitted exact-criterion premise in the first part follows, for
example, from consistency together with local stochastic continuity of
$\ell_n$, or from a uniform law of large numbers together with continuity of
the population criterion.  It is not a consequence of one-sided separation
alone and remains an explicit premise.

The fitted-gap conclusion is in-sample.  Without encoder complexity or
stability control, a network can fit training posteriors and behave
arbitrarily at unseen inputs.  Qualitative mixture closure can select a
diagonal sequence of ideal architectures at any prescribed gap order, but
this is existential undersmoothing: it gives neither an implementable growth
schedule nor a guarantee that stochastic optimization attains the required
$\epsilon_n$.

\subsubsection{A tractable special case: linear--Gaussian measurements}
\label{subsubsec:linear_gaussian_case}

We first give a case in which a fixed encoder architecture is sufficient.
Let $f_p(x^p)$, $f_c(x^c)$, and $f_s(x^s)$ be fixed finite-dimensional
feature maps and consider
\begin{equation}
\label{eq:lg_dgp_main}
\begin{aligned}
 z_i\mid x_i^p
 &\sim\mathcal N(Bf_p(x_i^p),\Sigma_0),\\
 w_{im}\mid z_i,x_i^c
 &\sim\mathcal N(a_m+D_mf_c(x_i^c)+A_mz_i,\Omega_m),\\
 y_i\mid z_i,x_i^s
 &\sim\mathcal N(a_y+D_yf_s(x_i^s)+Cz_i,\Omega_y).
\end{aligned}
\end{equation}
Here $\Omega_y$ is diagonal under the conditional independence assumption in
Eq.~\eqref{eq:dgp_deeplatent2}; allowing it to be unrestricted is a harmless
Gaussian relaxation.  The loading matrices $A_m$ and $C$ do not vary with
covariates.  This
restriction is what makes the exact posterior mean affine in the encoder
inputs.  Covariate-dependent loadings remain Gaussian, but their posterior
map involves matrix inversion as a function of the covariates and need not be
represented exactly by a finite generic neural network.

For a subset $S$, stack the observed measurements and outcomes as
\[
 u_{i,S}=a_S+D_Sf(x_i)+G_Sz_i+e_{i,S},
 \qquad e_{i,S}\sim\mathcal N(0,\Omega_S).
\]
Completing the square gives
\begin{equation}
\label{eq:lg_posterior_main}
\begin{aligned}
 V_S(\vartheta)
 &=\left(\Sigma_0^{-1}+G_S^\top\Omega_S^{-1}G_S\right)^{-1},\\
 m_S(\vartheta;V_i(S))
 &=V_S(\vartheta)\left[
 \Sigma_0^{-1}Bf_p(x_i^p)
 +G_S^\top\Omega_S^{-1}
 \{u_{i,S}-a_S-D_Sf(x_i)\}\right].
\end{aligned}
\end{equation}
Thus $\pi_\vartheta(\cdot\mid V_i(S))=\mathcal
N(m_S(\vartheta;V_i(S)),V_S(\vartheta))$.  We take $M$, $K$, the response
dimensions, and the feature dimensions to be fixed as $n$ grows.  There are
then only finitely many subsets, so an encoder with one affine mean head and
one unrestricted constant Cholesky head for each $S$ represents this
posterior exactly.  The architecture is fixed as $n$ grows; its weights are
estimated.

\begin{assumption}[Regular linear--Gaussian criterion]
\label{ass:lg_regular}
The unrestricted model is invariant under affine changes of latent
coordinates and may have additional confounding between prevalence and
measurement covariates.  Impose a complete model-specific location, scale,
shear, rotation, sign, and ordering normalization.  For a constant-zero-mean
factor model, for example, one can set $\Sigma_0=I_L$ and require a
nonsingular $L\times L$ loading block to be lower triangular with positive
diagonal; the
conditional-mean model may also need exclusion or coefficient restrictions
that distinguish $Bf_p$ from $D_mf_c$ and $D_yf_s$.  Let $\vartheta_0$ be an
interior point of the resulting finite-dimensional parameter space
$\mathcal V$.  Assume: (i) $U_i$ are i.i.d. and all subset marginals are
correctly specified at $\vartheta_0$; (ii) $\mathcal V$ is compact,
covariance eigenvalues are uniformly bounded away
from zero and infinity, and the composite criterion is globally separated at
$\vartheta_0$; (iii) the subset patterns with positive weight identify
$\vartheta_0$ under the normalization; (iv) $\bar\ell_\vartheta(U)$ is
continuous with an integrable global envelope and twice continuously
differentiable near $\vartheta_0$, with integrable score and Hessian envelopes;
and (v)
\[
 \E_0[\nabla_\vartheta\bar\ell_{\vartheta_0}(U)]=0,\qquad
 A_0:=-\E_0[\nabla^2_{\vartheta\vartheta}
 \bar\ell_{\vartheta_0}(U)],\qquad
 B_0:=\mathrm{Var}_0(\nabla_\vartheta\bar\ell_{\vartheta_0}(U))
\]
exist and $A_0$ is nonsingular.  Identification across subset patterns is a
substantive assumption: it can fail, for example, if no pattern links two
otherwise unaligned modalities.
\end{assumption}

\begin{theorem}[Exact profiling and Gaussian-case limit theory]
\label{thm:lg_profile_mle}
Let the variational family contain full-covariance Gaussians and let the
fixed encoder architecture contain the subset-specific affine heads described
above.  For every $\vartheta\in\mathcal V$,
\begin{equation}
 \sup_\phi\frac1n\sum_{i=1}^n
 \E_{S\mid O_i}[\mathcal L(\vartheta,q_\phi;V_i(S))]
 =\frac1n\sum_{i=1}^n\bar\ell_\vartheta(U_i).
 \label{eq:gaussian_profile_identity}
\end{equation}
Consequently, the generative component $\widehat\vartheta_n$ of any global
joint ELBO maximizer is the extremum estimator for the composite criterion in
Eq.~\eqref{eq:composite_subset_criterion}.  Under
Assumption~\ref{ass:lg_regular},
\[
 \widehat\vartheta_n\xrightarrow{p}\vartheta_0,
 \qquad
 \sqrt n(\widehat\vartheta_n-\vartheta_0)
 \rightsquigarrow\mathcal N(0,A_0^{-1}B_0A_0^{-1}).
\]
For a fixed subset this is an ordinary correctly specified likelihood, so
$A_0=B_0=I_0$ and the covariance reduces to $I_0^{-1}$.  The same likelihood
conclusion applies to the \emph{alternative} estimator based on one genuinely
observed random coarsening per unit and its correctly specified ignorable
design.  That alternative criterion is
$n^{-1}\sum_i\ell_\vartheta(V_i(S_i))$, not the expected-subset criterion in
Eq.~\eqref{eq:gaussian_profile_identity}.  Averaging several overlapping
subset log likelihoods is generally a composite likelihood, so the
information equality does not apply.
No identification or asymptotic-normality claim is required for the encoder
weights, which can be nonunique even when the induced conditional density is
unique.
\end{theorem}

The consistency argument is the standard extremum-estimator argument of
\citet{NeweyMcFadden1994}; the sandwich covariance is the corresponding
Godambe covariance for a composite criterion
\citep{VarinReidFirth2011}.

The same calculation proves exactness of corrected product-of-experts fusion
when every source appears once.  If outcomes form their own expert, the
fusion contains $|S|+1$ posterior experts and must divide by $|S|$ copies of
the prior.  The resulting Gaussian natural parameters give
Eq.~\eqref{eq:lg_posterior_main}.  The subset-head construction is an
existence benchmark; the corrected Gaussian PoE realizes the same posterior
with tied source-specific encoders.

\subsubsection{General posteriors: Gaussian-mixture and neural-network sieves}
\label{subsubsec:general_sieve}

We now keep the generative parameter finite-dimensional and allow the encoder
and its Gaussian-mixture output to grow with $n$.  This is a sieve estimator:
the decoder is the identified econometric model, while the encoder is an
increasingly rich nuisance function.  General results for sieve extremum
estimators are developed by \citet{Chen2007Sieve}; variational estimators as
profile $M$-estimators are studied by \citet{WestlingMcCormick2019}.

The next assumption gives primitive conditions under which ordinary finite
Gaussian mixtures are adequate for the reverse KL divergence used by the
ELBO.  The tail condition matters: density approximation in $L^1$, or in the
opposite KL direction, does not by itself imply
$\mathrm{KL}(q\|\pi)\to0$.  Related conditional-mixture approximation
results in the opposite KL direction are given by
\citet{Norets2010Conditional}; the proof below handles the ELBO direction
directly.

\begin{assumption}[Regular posterior class]
\label{ass:posterior_regular}
Let $\mathcal V\subset\mathbb R^{d_\vartheta}$ and the support
$\mathcal K_V\subset\mathbb R^{d_V}$ of a fixed finite-dimensional encoding
of $V=(S,u_S,x)$ be compact.  For
$(\vartheta,V)\in\mathcal V\times\mathcal K_V$, the posterior has a strictly
positive Lebesgue density $\pi_\vartheta(z\mid V)$ on $\mathbb R^L$ such
that: (i) $(\vartheta,V,z)\mapsto\pi_\vartheta(z\mid V)$ is continuous and
the family is equicontinuous in $z$ on compact sets; (ii) the family is
uniformly tight; and (iii), for constants $C<\infty$ and $r>0$,
\[
 \sup_{\vartheta,V}\pi_\vartheta(z\mid V)\le C,
 \qquad
 \sup_{\vartheta,V}\{-\log\pi_\vartheta(z\mid V)\}
 \le C(1+\|z\|^r).
\]
The polynomial can be replaced by an envelope integrable against the common
Gaussian tail envelopes used in the proof.
\end{assumption}

Gaussian priors combined with continuous linear decoders and Gaussian,
Bernoulli, categorical, or fixed-count multinomial likelihoods satisfy this
type of condition on compact parameter and data domains when covariance
eigenvalues and likelihood scales are uniformly nondegenerate.  For unbounded
data, the compact-domain statement can be replaced by an integrated version
plus a uniform-integrability condition.

\begin{proposition}[Constructive closure of both inference gaps]
\label{prop:mog_encoder_tightness}
Suppose Assumption~\ref{ass:posterior_regular} holds.  For every
$\varepsilon>0$ there is a finite number $J$, fixed Gaussian component
locations and covariances, and, for each $\vartheta$, continuous mixing-weight
functions $\alpha_{\vartheta,1}(V),\ldots,
\alpha_{\vartheta,J}(V)$ such that
\[
 \sup_{\vartheta\in\mathcal V}\sup_{V\in\mathcal K_V}
 \mathrm{KL}\!\left(
 \sum_{j=1}^J\alpha_{\vartheta,j}(V)
 \mathcal N(\mu_j,\Sigma_j)
 \,\middle\|\,\pi_\vartheta(\cdot\mid V)\right)<\varepsilon.
\]
Let the encoder density be
\[
 q_\phi(z\mid V)=\sum_{j=1}^J
 [\operatorname{softmax}\{g_\phi(V)\}]_j
 \mathcal N(z;\mu_j,\Sigma_j).
\]
If, for every finite $J$, the network sieve is dense in
$C(\mathcal K,\mathbb R^J)$ for every compact finite-dimensional input domain
$\mathcal K$ used here, and is closed under fixing auxiliary inputs (as are
standard fully connected ReLU networks with $J$ outputs and unbounded width),
then, possibly after increasing its size, there is for every $\vartheta$ an
encoder parameter $\phi(\vartheta)$ satisfying
\begin{equation}
 \sup_{\vartheta\in\mathcal V}\inf_\phi
 \sup_{V\in\mathcal K_V}
 \{\ell_\vartheta(V)-\mathcal L(\vartheta,q_\phi;V)\}
 <2\varepsilon.
 \label{eq:uniform_mog_tightness}
\end{equation}
Hence a Gaussian-mixture variational sieve closes the variational gap and a
deep encoder sieve closes the amortization gap.  This is an approximation
result; it does not assert that a finite sample or a particular optimizer
finds the approximating encoder.
\end{proposition}

The construction uses a common fine partition of a large latent-space cube.
The component locations are the cell centers and the mixing weights are
integrals of the posterior over the cells.  Joint continuity of the posterior
makes these weights continuous in $V$.  A neural network then approximates
their logits uniformly.  This direct construction supplies the continuity
that is needed for amortization.  The simple-hierarchical-model result of
\citet{MargossianBlei2024} separately ensures existence of an ideal inference
function in this model class, but does not by itself establish that a chosen
optimizer is continuous.

\subsubsection{Specialization to the GTM and affine IdealPointNN}
\label{subsubsec:model_specific_sieves}

For the implemented linear Generalized Topic Model, let
$H_K\in\mathbb R^{K\times(K-1)}$ satisfy
$H_K^\top H_K=I$ and $H_K^\top\mathbf1_K=0$, and write
\begin{equation}
 \eta_i\mid x_i^p\sim
 \mathcal N(\Gamma x_i^p+a,\Sigma),\qquad
 \theta_i=\operatorname{softmax}(H_K\eta_i),\qquad
 h_i=B\theta_i+D x_i^c+b.
 \label{eq:gtm_exact_linear_model}
\end{equation}
For word counts $w_i$ of length $N_i$, the cross-entropy branch is the
multinomial log likelihood, up to its data-only coefficient,
\[
 \log L_i(\eta)=w_i^\top h_i-N_i\log\sum_v\exp(h_{iv}).
\]
For an embedding $y_i$, the summed MSE branch
$-\|y_i-h_i\|^2$ is a Gaussian log likelihood with covariance
$\tfrac12I$, up to a parameter-free constant.  An MSE with another scaling is
likewise likelihood-calibrated only after declaring the corresponding fixed
Gaussian covariance.  A fixed finite collection of additional affine
Bernoulli, categorical, or Gaussian outcome heads can be included by adding
their normalized log likelihoods; the bounds below then use the sum of their
oscillations.

\begin{proposition}[GTM posterior is a bounded Gaussian tilt]
\label{prop:gtm_bounded_tilt}
Suppose the dimensions are fixed, the generative parameters range over a
compact set, the eigenvalues of $\Sigma$ are bounded away from zero and
infinity, and the covariates are bounded.  For cross-entropy also suppose
$N_i\leq\bar N$.  For cross-entropy, constants $0<c<C<\infty$, uniform over
the displayed sets, satisfy
\begin{equation}
 c\,\phi_{\Gamma x_i^p+a,\Sigma}(\eta)
 \leq\pi_\vartheta(\eta\mid w_i,x_i^p,x_i^c)
 \leq C\,\phi_{\Gamma x_i^p+a,\Sigma}(\eta).
 \label{eq:gtm_gaussian_comparison}
\end{equation}
The same uniform comparison holds for bounded $y_i$ in the MSE branch.
Cross-entropy and bounded-$y$ MSE posteriors are therefore positive and
continuous, have uniform Gaussian tails and moments, and satisfy the
regular-posterior conditions used by
Proposition~\ref{prop:mog_encoder_tightness}.  A bounded-perturbation argument
also gives a uniform log-Sobolev inequality and therefore a uniform $T_2$
transportation inequality for these two cases \citep{OttoVillani2000}.  For
unbounded Gaussian $y_i$, the two-sided comparison and perturbation constants
are record-dependent; under a suitable exponential envelope we claim only the
integrated tail and moment controls required below, not a uniform log-Sobolev
or $T_2$ constant.  The posterior need not be log-concave: the softmax
composition can create separated modes.  Strong log-concavity requires the
additional prior-dominance condition that the smallest prior precision
dominate the likelihood's worst negative curvature.
\end{proposition}

It follows that, under exact-likelihood identification and curvature,
Theorem~\ref{thm:one_comparator} applies to the GTM.  In particular, gap plus
optimization order $o_p(n^{-1})$ gives raw-ELBO MLE equivalence, while order
$o_p(n^{-1/2})$ supplies the preliminary rate for the importance correction
below.  Qualitative closure supplies those rates only through an existential
diagonal architecture whose mixture bandwidths can shrink and whose input
contains every posterior-relevant variable.

For the affine IdealPointNN, let
\[
 z_i\mid x_i^p\sim\mathcal N(\Gamma x_i^p+a,\Sigma)
\]
and suppose each selected measurement or outcome head has an affine natural
parameter in $z$.  Allow Bernoulli logits, categorical or fixed-count
multinomial logits, and Gaussian means with fixed nondegenerate covariance.

\begin{proposition}[Affine IdealPoint posterior shape]
\label{prop:idealpoint_strong_logconcavity}
On a compact local parameter set with
$\lambda_{\max}(\Sigma)\leq\bar\lambda$, the posterior is uniformly strongly
log-concave.  In schematic block notation its negative log Hessian is
\begin{align}
 \nabla_z^2\{-\log\pi_\vartheta(z\mid U)\}
 &=\Sigma^{-1}
 +\sum_g A_g^\top\Omega_g^{-1}A_g
 +\sum_b A_b^\top\operatorname{diag}\{p_b(1-p_b)\}A_b\\
 &\quad+\sum_c n_c A_c^\top
   (\operatorname{Diag}p_c-p_cp_c^\top)A_c
 \succeq\bar\lambda^{-1}I.
 \label{eq:idealpoint_posterior_hessian}
\end{align}
Thus it is positive, unimodal, Gaussian-tailed, and satisfies uniform
log-Sobolev and $T_2$ inequalities.  Bernoulli and categorical marginal
likelihoods remain analytically intractable.  With affine Gaussian heads only,
the posterior is exactly Gaussian.  Hidden nonlinear decoder/predictor layers,
image decoders, and unnormalized losses are outside this proposition.
\end{proposition}

Both specializations require an identified chart.  For GTM, topic labels must
be anchored or the reported functional permutation invariant; multinomial
output logits need a reference row or sum-zero restriction; only one of a
prior bias and a prevalence intercept may be free; only one of a decoder bias
and a content intercept may be free; and, after choosing that baseline, a
constraint such as $B\mathbf1_K=0$ removes the decoder intercept gauge.  Full
rank designs, nonsingular information, and exact-criterion separation remain
substantive assumptions.  The contrast map removes the latent softmax shift,
but it does not remove topic permutations or prove global identification.
Scaling $\eta$ is not in general an exact gauge for a decoder affine in
$\theta$.

For IdealPointNN with learned $\Sigma$ and unrestricted loadings, every
nonsingular affine change of latent coordinates can be absorbed by the prior,
loadings, and intercepts.  Fixing $\Sigma=I$ still leaves rotations and
reflections.  Coordinate-level inference therefore needs, for example, a
single location convention, fixed covariance (or a quotient-whitened chart),
a positive discrimination anchor in one dimension or a lower-triangular
positive-diagonal anchor block in higher dimensions, categorical-logit gauges,
and exclusions where the same covariates enter both the latent mean and direct
outcome effects.  Otherwise only smooth invariant functionals are covered.

\begin{proposition}[Omitted conditioning creates an irreducible gap]
\label{prop:omitted_conditioning_gap}
Let $U$ be everything supplied to an encoder and let an omitted discrete
variable $C$, conditional on $U=u$, take values $c$ with probabilities
$\omega_c>0$.  Write
$\pi_c(z)=\pi_\vartheta(z\mid U=u,C=c)$.  Even over all densities depending
on $u$ alone,
\begin{equation}
 \inf_q\sum_c\omega_c\mathrm{KL}(q\|\pi_c)
 =-\log\int\exp\!\left\{\sum_c\omega_c\log\pi_c(z)\right\}\diff z.
 \label{eq:omitted_conditioning_gap}
\end{equation}
The optimizer is the normalized weighted geometric mean of the posteriors.
The value is strictly positive unless all positive-weight posteriors agree
almost everywhere.
\end{proposition}

This is a capacity-independent obstruction.  Content and prediction
covariates that affect a likelihood must be encoder inputs.  Likewise, an
outcome included in the joint likelihood must enter the proposal exactly
once: appending it to each modality expert before multiplying experts counts
its evidence repeatedly, whereas omitting it targets a different posterior.

For comparison, we retain a stronger traditional joint-sieve theorem that
also gives an encoder rate.  It is not needed for the one-comparator argument
and its entropy, curvature, and stochastic-modulus premises must be verified
for the chosen network class.  Let
\[
 m(\vartheta,q;V):=\mathcal L(\vartheta,q;V),\qquad
 \bar m(\vartheta,q;U):=\E_{S\mid O}[m(\vartheta,q;V(S))],
 \qquad
 \bar M(\vartheta,q):=\E_0[\bar m(\vartheta,q;U)],
\]
and let $\mathcal Q_n$ contain $J_n$-component mixtures whose parameters are
outputs of a bounded ReLU network with $P_n$ weights.  Write
$\mathcal H_n=\mathcal V\times\mathcal Q_n$ and let $\kappa_n$ be an entropy
index for this class; for bounded finite networks, the canonical order is
$\kappa_n\asymp P_n\log P_n$.  Let $\rho$ be a metric satisfying locally
\begin{equation}
 \rho^2((\vartheta,q),(\vartheta_0,\pi_0))
 \asymp \|\vartheta-\vartheta_0\|^2
 +\E_0\E_{S\mid O} h^2\!\left(q(\cdot\mid V(S)),
 \pi_\vartheta(\cdot\mid V(S))\right),
 \label{eq:sieve_metric}
\end{equation}
where $h$ is Hellinger distance and
$\pi_0=\pi_{\vartheta_0}$.

\begin{assumption}[Sieve estimation]
\label{ass:sieve_estimation}
After an identifying normalization, assume:
\begin{enumerate}
\item[(i)] $U_i$ are i.i.d., the subset marginals are correctly specified at
$\vartheta_0$, and $\vartheta_0$ uniquely maximizes
$\E_0\bar\ell_\vartheta(U)$.  Thus the collection of subset patterns assigned
positive probability identifies the normalized generative parameter.
\item[(ii)] (Quadratic curvature) In a neighborhood of
$(\vartheta_0,\pi_0)$,
\[
 \bar M(\vartheta_0,\pi_0)-\bar M(\vartheta,q)
 \ge c\rho^2((\vartheta,q),(\vartheta_0,\pi_0)).
\]
\item[(iii)] (Sieve approximation) There is
$(\vartheta_0,q_n^*)\in\mathcal H_n$ with
$\rho((\vartheta_0,q_n^*),(\vartheta_0,\pi_0))\le a_n$ and population
criterion loss $O(a_n^2)$, where $a_n\to0$.
Define $r_n:=a_n+\sqrt{\kappa_n/n}$.
\item[(iv)] (Local stochastic modulus) For some sufficiently small fixed
$\bar\delta>0$,
\[
 \sup_{\delta\in[r_n,\bar\delta]}
 \frac{
 \sup_{\rho((\vartheta,q),(\vartheta_0,q_n^*))\le\delta}
 \left|\sqrt n(P_n-P_0)
 \{\bar m(\vartheta,q)-\bar m(\vartheta_0,q_n^*)\}\right|}
 {\delta\sqrt{\kappa_n}}=O_p(1).
\]
\item[(v)] (Global control) $\kappa_n/n\to0$,
$\sup_{\mathcal H_n}|(P_n-P_0)\bar m|=o_p(1)$, and for every
$\epsilon>0$,
\[
 \liminf_n\inf_{\substack{(\vartheta,q)\in\mathcal H_n:\\
 \rho((\vartheta,q),(\vartheta_0,\pi_0))\ge\epsilon}}
 \{\bar M(\vartheta_0,\pi_0)-\bar M(\vartheta,q)\}>0.
\]
The computed estimator satisfies
\[
 P_n\bar m(\widehat\vartheta_n,\widehat q_n)
 \ge\sup_{(\vartheta,q)\in\mathcal H_n}P_n\bar m(\vartheta,q)-o_p(r_n^2).
\]
\end{enumerate}
\end{assumption}

\begin{theorem}[Consistency and rate of the general ELBO estimator]
\label{thm:sieve_consistency_rate}
Under Assumptions~\ref{ass:posterior_regular} and
\ref{ass:sieve_estimation},
\[
 \rho((\widehat\vartheta_n,\widehat q_n),
 (\vartheta_0,\pi_0))=O_p(r_n),
 \qquad
 \|\widehat\vartheta_n-\vartheta_0\|=O_p(r_n).
\]
In particular, the generative parameter and the encoder-induced posterior are
consistent.  If a known transformation-group ambiguity remains, the same
conclusion holds in the corresponding quotient metric.
\end{theorem}

The theorem makes the rate transparent: it is the sum of sieve approximation
error and sampling error.  The exact value of $a_n$ must be verified for the
chosen posterior class and parameter bounds.  For illustration, suppose a
particular Gaussian-mixture construction and ReLU sieve have been shown to
satisfy
\begin{equation}
 a_n\lesssim J_n^{-\beta/L}
 + (P_n/J_n)^{-\alpha/d_V},
 \qquad
 r_n\lesssim J_n^{-\beta/L}
 +(P_n/J_n)^{-\alpha/d_V}
 +\sqrt{P_n\log P_n/n}.
 \label{eq:illustrative_sieve_rate}
\end{equation}
Such bounds require control of shrinking mixture variances, component weights
near the simplex boundary, and network norms; they do not follow from
qualitative universal approximation.  Gaussian-mixture and ReLU approximation
results such as \citet{KruijerRousseauVdV2010} and
\citet{Yarotsky2018Deep} provide tools for verifying model-specific versions.
Equation~\eqref{eq:illustrative_sieve_rate} is not a dimension-free parametric
rate for the encoder.  It displays the curse of dimensionality in both latent
space and encoder-input space.  The abstract rate $a_n$ in
Theorem~\ref{thm:sieve_consistency_rate} should be used when a different
mixture construction or network approximation bound is available.

Ignoring logarithms, put $J_n=n^j$ and $P_n=n^p$.  Balancing the three terms
in Eq.~\eqref{eq:illustrative_sieve_rate} gives the illustrative achievable
order
 \[
 \gamma=\frac{1}{2+L/\beta+d_V/\alpha},\qquad
 r_n=O(n^{-\gamma})\quad\text{up to logarithmic factors}.
\]
All three terms can be $o(n^{-1/4})$ only if
$L/\beta+d_V/\alpha<2$.  This joint restriction, rather than a condition on
the encoder smoothness alone, is the relevant implication of the displayed
rate.

Root-$n$ inference for $\vartheta$ requires more than consistency of the
encoder.  The unrestricted posterior curve reproduces the exact likelihood,
but this level identity does not make the score of a restricted ELBO
orthogonal to encoder error.  For one record write
\[
 v(\vartheta,q;V)
 =\E_q\{\log p_\vartheta(u,Z\mid x,S)-\log q(Z\mid V)\}.
\]
Holding $q$ fixed gives
$\nabla_\vartheta v=\E_q s_c(\vartheta;V,Z)$, where $s_c$ is the complete-data
score; write $s_o(\vartheta;V)=\nabla_\vartheta\ell_\vartheta(V)$ for the
observed score.

\begin{proposition}[No automatic ELBO-score orthogonality]
\label{prop:no_automatic_orthogonality}
Let
$q_t(z\mid V)=\pi_\vartheta(z\mid V)\{1+t b(z,V)\}$ with
$\E_{\pi_\vartheta}(b\mid V)=0$.  Then
\begin{equation}
 \left.\frac{\diff}{\diff t}\nabla_\vartheta
 v(\vartheta,q_t;V)\right|_{t=0}
 =\E_{\pi_\vartheta}
 \left[b(Z,V)\{s_c(\vartheta;V,Z)-s_o(\vartheta;V)\}\mid V\right],
 \label{eq:elbo_score_path_derivative}
\end{equation}
which is generally nonzero.  The ELBO \emph{level} is stationary with respect
to $q$ at the posterior; its generative score generally is not.  Under
posterior exponential moments, reverse KL alone gives the worst-case order
\[
 \|\E_qs_c-\E_\pi s_c\|
 \lesssim \mathrm{KL}(q\|\pi)^{1/2}+\mathrm{KL}(q\|\pi),
\]
and the square-root term is locally sharp.  Hence an $o_p(n^{-1/4})$ nuisance
rate does not, by itself, establish root-$n$ inference for the raw ELBO
estimator.

For a Bernoulli decoder with intercept score $s_{c,\alpha}$, take
$b=c_V\{s_{c,\alpha}-s_{o,\alpha}\}$, where $c_V>0$ is small enough that the
density path is nonnegative.  The $\alpha$ coordinate of
Eq.~\eqref{eq:elbo_score_path_derivative} is then
$c_V\operatorname{Var}_{\pi_\vartheta}(s_{c,\alpha}\mid V)$, which is strictly
positive whenever the latent variable changes the Bernoulli probability on a
nondegenerate posterior support.
\end{proposition}

Indeed, if
$R_s(\vartheta)=\inf_{q\in\mathcal Q_s}
\E_0\mathrm{KL}(q\|\pi_\vartheta)$ has a regular minimizing selector, the
envelope theorem gives
\[
 \nabla_\vartheta R_s(\vartheta)
 =\E_0\{s_o(\vartheta;V)-\E_{q_{s,\vartheta}}s_c(\vartheta;V,Z)\}.
\]
Thus confidence-interval centering depends on the derivative of the
variational gap, not only its level.  A specially structured mixture
projection could have tangent cancellation, but that requires a separate
projection theorem; universal approximation does not supply it.

For completeness, the following high-level profile conditions preserve a
valid but explicitly assumed route for the raw ELBO estimator.  Let
\[
 \widehat M_n(\vartheta):=\sup_{q\in\mathcal Q_n}
 P_n\bar m(\vartheta,q;U).
\]

\begin{assumption}[Profile expansions]
\label{ass:profile_clt}
Work in an identified local chart.  Let
\[
 \bar s_0(U):=\nabla_\vartheta\bar\ell_{\vartheta_0}(U),\qquad
 A_0:=-\E_0\nabla^2_{\vartheta\vartheta}
 \bar\ell_{\vartheta_0}(U),\qquad
 B_0:=\mathrm{Var}_0(\bar s_0(U)).
\]
Assume $A_0$ is nonsingular, $\E_0\bar s_0=0$,
$\E_0\|\bar s_0\|^2<\infty$, and the empirical profile has a twice
differentiable local version.  For some deterministic
$\delta_n\downarrow0$ such that
$\Pr(\|\widehat\vartheta_n-\vartheta_0\|\le\delta_n)\to1$, assume
\begin{align}
 \nabla_\vartheta\widehat M_n(\vartheta_0)
 &=P_n\bar s_0+o_p(n^{-1/2}),
 \label{eq:assumed_profile_score}\\
 \sup_{\|\vartheta-\vartheta_0\|\le\delta_n}
 \|\nabla^2_{\vartheta\vartheta}\widehat M_n(\vartheta)+A_0\|
 &=o_p(1).
 \label{eq:assumed_profile_hessian}
\end{align}
\end{assumption}

A possible route to
Eqs.~\eqref{eq:assumed_profile_score}--\eqref{eq:assumed_profile_hessian} is
to work in a nuisance norm in which the ELBO is uniformly three-times
Fr\'echet differentiable, obtain $r_n=o(n^{-1/4})$ in that norm, impose
$o_p(n^{-1/2})$ nuisance first-order conditions, approximate each
least-favorable tangent
$\partial\pi_\vartheta/\partial\vartheta_j|_{\vartheta_0}$ by a sieve tangent
with induced score bias $o(n^{-1/2})$, and verify stochastic equicontinuity of
the derivative classes.  The tangent-bias condition is the needed first-order
cancellation; the $n^{-1/4}$ rate only controls the remaining quadratic term.
These are model-specific undersmoothing conditions from sieve profile
$M$-estimation, not consequences of posterior stationarity or universal
approximation alone \citep{Chen2007Sieve,WestlingMcCormick2019}.

\begin{theorem}[Root-$n$ rate and asymptotic normality in the general case]
\label{thm:sieve_clt}
Under Assumptions~\ref{ass:posterior_regular},
\ref{ass:sieve_estimation}, and \ref{ass:profile_clt}, if
$\widehat\vartheta_n$ is an interior approximate maximizer whose profile-score
residual is $o_p(n^{-1/2})$, then
\[
 \sqrt n(\widehat\vartheta_n-\vartheta_0)
 =A_0^{-1}\frac1{\sqrt n}\sum_{i=1}^n \bar s_0(U_i)+o_p(1)
 \rightsquigarrow
 \mathcal N(0,A_0^{-1}B_0A_0^{-1}).
\]
For the random-subset composite criterion, $A_0$ and $B_0$ need not coincide.
For an ordinary correctly specified observed-data likelihood, the information
identity gives $A_0=B_0=I_0$.  On a quotient space the statement applies to
any smooth identified representative or smooth invariant functional.
\end{theorem}

Theorem~\ref{thm:sieve_clt} therefore does not claim automatic Neyman
orthogonality.  Its profile-score expansion directly assumes away any
first-order variational bias.  The transparent alternative
$\rho_n+\epsilon_n=o_p(n^{-1})$ in
Theorem~\ref{thm:one_comparator} avoids that assumption by making the entire
raw estimator MLE-equivalent.  The next construction instead removes the
first-order score bias by evaluating the exact likelihood score with defensive
importance sampling.

\subsubsection{Defensive likelihood correction and confidence intervals}
\label{subsubsec:defensive_one_step}

Choose one exact target.  For the full-record likelihood, let
$s_i^\star(\vartheta)=\nabla_\vartheta\log p_\vartheta(U_i)$; for the
expected-subset composite criterion, let
\[
 s_i^\star(\vartheta)
 =\E_{S\mid O_i}\nabla_\vartheta\ell_\vartheta(V_i(S)).
\]
Write $S_n^\star=P_ns_i^\star$,
$A_0=-\E_0\nabla_\vartheta s_i^\star(\vartheta_0)$, and
$B_0=\operatorname{Var}_0\{s_i^\star(\vartheta_0)\}$.  These two targets need
not have the same influence function; for the composite target, each subset
posterior ratio below must be normalized separately before averaging over
$S$.

Freeze the trained encoder $\widehat q_{i,S}$ and, at the parameter value at
which a score is evaluated, use the defensive proposal
\begin{equation}
 g_{i,S,\vartheta}(z)
 =(1-\delta)\widehat q_{i,S}(z)
  +\delta p_\vartheta(z\mid x_i^p),\qquad0<\delta\leq1.
 \label{eq:defensive_proposal}
\end{equation}
With
\[
 f_{i,S,\vartheta}(z)
 =p_\vartheta(z\mid x_i^p)L_{i,S,\vartheta}(z),\qquad
 W_{i,S,\vartheta}(z)=f_{i,S,\vartheta}(z)/g_{i,S,\vartheta}(z),
\]
every integrable complete-data quantity obeys the exact identity
\begin{equation}
 \E_{\pi_{i,S,\vartheta}}h(Z)
 =\frac{\E_g\{W_{i,S,\vartheta}(Z)h(Z)\}}
        {\E_g W_{i,S,\vartheta}(Z)}.
 \label{eq:defensive_posterior_identity}
\end{equation}
Taking $h=s_c$ gives Fisher's identity.  If
$H_c=\nabla_{\vartheta\vartheta}^2\log p_\vartheta(U_i,Z)$, Louis' identity
gives the observed information contribution
\begin{equation}
 I_{o,i}(\vartheta)
 =\E_\pi(-H_c\mid U_i)-\operatorname{Var}_\pi(s_c\mid U_i)
 \label{eq:louis_identity_defensive}
\end{equation}
\citep{Louis1982}.  The proposal has support regardless of encoder quality and
$W\leq\delta^{-1}L$.  It must use the prior at the current score-evaluation
parameter; Gaussian priors at two different parameter values need not have a
bounded ratio.  The proposal and samples are held fixed when differentiating
the target joint density.

For $R$ fresh draws, use the self-normalized ratio estimator
\[
 \widehat T_{i,R}(h)
 =\frac{\sum_{r=1}^R W_{ir}h(Z_{ir})}{\sum_{r=1}^R W_{ir}}.
\]
Conditionally on the complete training sample and the frozen encoder, suppose
its record-level error $e_{i,R}$ satisfies
\[
 \|\E(e_{i,R}\mid\mathcal F_n)\|\leq b_i/R,
 \qquad
 \E\{\|e_{i,R}-\E(e_{i,R}\mid\mathcal F_n)\|^2\mid\mathcal F_n\}
 \leq v_i/R,
\]
with independent streams across records.  Then
\begin{equation}
 P_ne_{i,R}
 =O_p\!\left(\frac{\bar b_n}{R}
       +\sqrt{\frac{\bar v_n}{nR}}\right).
 \label{eq:snis_average_rate}
\end{equation}
Under bounded averaged envelopes, raw self-normalization therefore requires
\begin{equation}
 R_n/\sqrt n\longrightarrow\infty
 \label{eq:raw_snis_rate}
\end{equation}
for root-$n$-negligible score error.  The ratio bias of order $R^{-1}$ does not
average away.  A bounded normalized posterior/proposal ratio, adequate
complete-score moments, and an inverse-denominator moment bound are sufficient
primitive conditions.  Bounded unnormalized weights alone are not.

\begin{theorem}[Defensive one-step correction]
\label{thm:defensive_one_step}
Suppose $\widehat\vartheta_n$ is in the identified interior and
$\|\widehat\vartheta_n-\vartheta_0\|=o_p(n^{-1/4})$.  Assume the exact target
score satisfies the uniform local expansion
\begin{equation}
 S_n^\star(\vartheta)
 =S_n^\star(\vartheta_0)-A_0(\vartheta-\vartheta_0)
 +o_p(n^{-1/2})+O_p(\|\vartheta-\vartheta_0\|^2),
 \label{eq:exact_score_local_expansion}
\end{equation}
$A_0$ is nonsingular, and the score CLT holds.  At the random preliminary
estimate, suppose the defensive estimates obey
\begin{align}
 \widehat S_{n,R}(\widehat\vartheta_n)
  -S_n^\star(\widehat\vartheta_n)&=o_p(n^{-1/2}),
 \label{eq:defensive_score_accuracy}\\
 \widehat A_{n,R}(\widehat\vartheta_n)&\to_p A_0,
 \label{eq:defensive_info_consistency}\\
 \|\{\widehat A_{n,R}(\widehat\vartheta_n)-A_0\}
 (\widehat\vartheta_n-\vartheta_0)\|&=o_p(n^{-1/2}).
 \label{eq:critical_information_product}
\end{align}
The last product condition is essential; consistency of $\widehat A$ alone is
not sufficient for a one-step update from an $o_p(n^{-1/4})$ preliminary
estimate.  Suppose also that correction streams are fresh and independent
across records, Eq.~\eqref{eq:snis_average_rate} holds at the random evaluation
point, and any outer subset Monte Carlo error is $o_p(n^{-1/2})$.  Then
\begin{equation}
 \check\vartheta_n
 =\widehat\vartheta_n
  +\widehat A_{n,R}(\widehat\vartheta_n)^{-1}
   \widehat S_{n,R}(\widehat\vartheta_n)
 \label{eq:defensive_one_step_update}
\end{equation}
satisfies
\[
 \sqrt n(\check\vartheta_n-\vartheta_0)
 =A_0^{-1}\frac1{\sqrt n}\sum_{i=1}^n
 s_i^\star(\vartheta_0)+o_p(1)
 \rightsquigarrow\mathcal N(0,A_0^{-1}B_0A_0^{-1}).
\]
The preliminary-rate premise follows from
Theorem~\ref{thm:one_comparator} if
$\rho_n+\epsilon_n=o_p(n^{-1/2})$ and the exact estimator is root-$n$
regular.
\end{theorem}

\begin{theorem}[Fixed unbiased randomized-score corollary]
\label{thm:fixed_unbiased_randomized_score}
Let a randomized unit score be
\[
 \psi_i(\vartheta)=s_i^\star(\vartheta)+\xi_i(\vartheta),
 \qquad
 \E\{\xi_i(\vartheta_0)\mid U_i\}=0,
\]
with simulation streams independent across units.  Suppose the conditional
triangular-array CLT holds for the randomized unit scores and, at the
preliminary estimator, local stochastic continuity gives
\[
 P_n\{\xi_i(\widehat\vartheta_n)-\xi_i(\vartheta_0)\}
 =o_p(n^{-1/2}).
\]
If the exact-score expansion and preliminary-rate premises of
Theorem~\ref{thm:defensive_one_step} hold, and the information estimator is
consistent for $A_0$ and obeys the product condition
Eq.~\eqref{eq:critical_information_product}, then the one-step estimator based
on $P_n\psi_i(\widehat\vartheta_n)$ satisfies
\[
 \sqrt n(\check\vartheta_n-\vartheta_0)
 =A_0^{-1}\frac1{\sqrt n}\sum_{i=1}^n
 \{s_i^\star(\vartheta_0)+\xi_i(\vartheta_0)\}+o_p(1),
\]
with covariance
\begin{equation}
 A_0^{-1}(B_0+C_{\rm MC})A_0^{-1},\qquad
 C_{\rm MC}:=\E\{\operatorname{Var}(\xi_i\mid U_i)\}.
 \label{eq:fixed_unbiased_mc_covariance}
\end{equation}
For $K$ exact posterior draws used to average the complete-data Fisher score
for a full record,
\begin{equation}
 C_{\rm MC}=K^{-1}\E\{\operatorname{Var}_{\pi_i}
 (s_{c,i}\mid U_i)\}.
 \label{eq:exact_posterior_mc_covariance}
\end{equation}
For $K_S$ independent exact subset-score draws used to average the
expected-subset score,
\begin{equation}
 C_{\rm MC}=K_S^{-1}\E\{\operatorname{Var}_{S}
 (s_{i,S}\mid U_i)\}.
 \label{eq:exact_subset_mc_covariance}
\end{equation}
If posterior and subset simulation are both present, define $\xi_i$ as their
aggregate error in Eq.~\eqref{eq:fixed_unbiased_mc_covariance}.  Its conditional
variance includes both cross-covariance matrices dictated by the simulation
design.  When the two errors are conditionally independent, their conditional
variance contributions add.  This corollary applies only to an unbiased
randomized score; it explicitly excludes biased fixed-$R$ self-normalized
importance ratios.
\end{theorem}

A local Lipschitz law for the exact observed information, together with
root-$n$ stochastic fluctuation and sufficiently accurate importance moments,
is one primitive route to
$\widehat A-A_0=O_p(\|\widehat\vartheta-\vartheta_0\|+n^{-1/2})$ and hence
Eq.~\eqref{eq:critical_information_product}.  For intervals, recompute
$\widehat A$ at $\check\vartheta_n$ with fresh draws.  If estimated exact unit
scores $\widehat s_i^\star$ are used to estimate $B_0$, require the stronger
unit-level condition
\begin{equation}
 P_n\|\widehat s_i^\star-s_i^\star\|^2=o_p(1).
 \label{eq:unit_score_l2_covariance_condition}
\end{equation}
Then their empirical covariance consistently estimates $B_0$.  With a fixed
unbiased noisy unit score from
Theorem~\ref{thm:fixed_unbiased_randomized_score}, its empirical covariance
instead targets $B_0+C_{\rm MC}$.  For an ordinary correctly specified full
likelihood, $A_0=B_0$; for the expected-subset target, $B_0$ is the covariance
of the \emph{subset-averaged score}, not the average of subset score outer
products.

The defensive conditions are especially transparent here.  For bounded
document length, GTM cross-entropy has a uniformly bounded likelihood
oscillation and therefore a bounded normalized ratio.  Affine IdealPoint
Bernoulli/categorical models with a fixed number of responses and compact
parameters have a positive lower bound for the finite-pattern marginal
probabilities.  GTM MSE needs integrated residual envelopes.  Pure affine
IdealPoint MSE should use its analytic Gaussian posterior; in a mixed model,
the Gaussian-channel posterior can be the defensive base and only the bounded
discrete likelihood need be importance corrected.  Growing document or item
counts make the envelopes grow and must be reflected in $R_n$.

\begin{theorem}[Iterated defensive correction]
\label{thm:iterated_defensive_correction}
Let $\widetilde\vartheta_n$ solve $S_n^\star(\vartheta)=0$, let
$A_n(\vartheta)=-\nabla_\vartheta S_n^\star(\vartheta)$, and suppose on a
fixed basin containing $\widetilde\vartheta_n$ that $A_n$ is uniformly
nonsingular, bounded, and Lipschitz.  Then the exact Newton map
$\mathcal N_n(\vartheta)=\vartheta+A_n(\vartheta)^{-1}S_n^\star(\vartheta)$
obeys
\[
 \|\mathcal N_n(\vartheta)-\widetilde\vartheta_n\|
 \leq C\|\vartheta-\widetilde\vartheta_n\|^2.
\]
Starting from a consistent ELBO estimator inside a radius $r$ with
$q:=2Cr<1$, use fresh defensive draws at each iteration.  If the approximate
Newton maps have maximal error
\[
 \max_{k<K_n}\|\widehat{\mathcal N}_{n,k}(\vartheta_n^{(k)})
 -\mathcal N_n(\vartheta_n^{(k)})\|\leq\xi_n,
 \qquad \xi_n=o_p(n^{-1/2}),
\]
and all iterates stay in the basin, then, writing
$e_{n,0}=\|\vartheta_n^{(0)}-\widetilde\vartheta_n\|$,
\begin{equation}
 \|\vartheta_n^{(K_n)}-\widetilde\vartheta_n\|
 \leq C^{-1}(Ce_{n,0})^{2^{K_n}}+\frac{\xi_n}{1-q}.
 \label{eq:iterated_newton_bound}
\end{equation}
Thus if the first term is $o_p(n^{-1/2})$, the iterated estimator is
root-$n$ equivalent to the exact estimator.  A growing or data-dependent
$K_n$ requires a maximal inequality or union bound for every attempted
importance stream; freshness alone does not prove the displayed maximum.
\end{theorem}

If $e_{n,0}=O_p(n^{-a})$ for any $a>0$, a fixed number of iterations with
$2^{K}a>1/2$ suffices.  With an unknown slower rate, a stopping rule based on
an independently accurate exact-score residual can be used, but the threshold
must be $o(n^{-1/2})$ and importance error must be controlled uniformly over
all inspected iterates.  A raw ELBO gradient is not such a residual.

\subsubsection{Louis moments and what reverse KL does not control}
\label{subsubsec:louis_moments}

Reverse KL plus Gaussian posterior tails controls expectations of quantities
with at most quadratic growth under suitable exponential-moment bounds; it
does not control unrestricted fourth moments.  To see this, take
$P=\mathcal N(0,1)$ and
\[
 Q_\delta=(1-2\delta^3)P+2\delta^3\mathcal N(\delta^{-1},1).
\]
Convexity gives
$\mathrm{KL}(Q_\delta\|P)\leq\delta$, while
$\E_{Q_\delta}Z^4=3+2\delta^{-1}+O(\delta)\to\infty$.  Thus fitted reverse KL
alone does not justify untruncated Louis terms involving learned covariance
and loading scores.

One remedy is an integrated uniform-integrability or common tail condition on
the fitted sieve.  A second is componentwise truncation
$T_\tau(g)=\max(-\tau,\min(g,\tau))$.  If
$\bar\delta_n=P_n\mathrm{KL}(\widehat q_i\|\pi_i)$, Pinsker's inequality and
Jensen's inequality give
\begin{equation}
 P_n|\E_{\widehat q_i}T_{\tau_n}(g_i)
       -\E_{\pi_i}T_{\tau_n}(g_i)|
 \leq\sqrt2\,\tau_n\bar\delta_n^{1/2}.
 \label{eq:truncated_louis_bound}
\end{equation}
If $\bar\delta_n=o_p(n^{-1})$, choose
$\tau_n\to\infty$ and $\tau_n=o(\sqrt n)$; exact-posterior uniform
integrability removes the target truncation bias and
Eq.~\eqref{eq:truncated_louis_bound} vanishes.  Defensive importance-weighted
Louis estimation instead controls these polynomial moments through the
proposal envelopes and does not rely on fitted-encoder tail behavior.

\begin{proposition}[Posterior recovery and the limit of latent-factor claims]
\label{prop:posterior_recovery}
Under Theorem~\ref{thm:sieve_consistency_rate},
\[
 \E_0\E_{S\mid O} h^2\!\left(\widehat q_n(\cdot\mid V(S)),
 \pi_{\widehat\vartheta_n}(\cdot\mid V(S))\right)\xrightarrow{p}0.
\]
If the posterior family is continuous in $\vartheta$ in integrated Hellinger
distance, then $\widehat q_n$ also converges to $\pi_{\vartheta_0}$.  Posterior
functionals of bounded $g$ converge in the integrated sense
\[
 \E_0\E_{S\mid O}\left|
 \E_{\widehat q_n}[g(z)]-\E_{\pi_{\vartheta_0}}[g(z)]\right|
 \xrightarrow{p}0.
\]
For an unbounded $g$, the same conclusion holds under the explicit integrated
uniform-integrability conditions that, for every $\epsilon>0$,
\begin{align*}
 \lim_{M\to\infty}\limsup_{n\to\infty}
 \Pr\!\left(\E_0\E_{S\mid O}
 \E_{\widehat q_n}[|g|\mathbf 1\{|g|>M\}]>\epsilon\right)&=0,\\
 \lim_{M\to\infty}\E_0\E_{S\mid O}
 \E_{\pi_{\vartheta_0}}[|g|\mathbf 1\{|g|>M\}]&=0.
\end{align*}
These are integrated, not pointwise-in-$V$, conclusions.  If
the integrated exact posterior variance remains bounded away from zero with a
fixed amount of information per unit, no estimator can have vanishing
mean-squared error for the realized $z_i$; under posterior
anti-concentration, convergence in probability also fails.  Thus the results
do not imply $\widehat z_i\to z_i$.  Point recovery needs a within-unit
asymptotic regime under which the exact posterior itself concentrates.
\end{proposition}

\subsubsection{Optimization and Monte Carlo error}

The preceding theorems concern the exact expected-subset ELBO and a
sufficiently accurate extremum estimator.  Let
\[
 C_n(\vartheta,q):=P_n\bar m(\vartheta,q),\qquad
 \widetilde C_n(\vartheta,q):=C_n(\vartheta,q)+R_n(\vartheta,q).
\]
For first-order inference define the \emph{profiled} perturbation
\[
 R_n^p(\vartheta):=
 \sup_{q\in\mathcal Q_n}\widetilde C_n(\vartheta,q)
 -\sup_{q\in\mathcal Q_n}C_n(\vartheta,q).
\]
Let $\widetilde\eta_n=(\widetilde\vartheta_n,\widetilde q_n)$ satisfy
\[
 \widetilde C_n(\widetilde\eta_n)
 \ge \sup_{\eta\in\mathcal H_n}\widetilde C_n(\eta)-\epsilon_n,
 \qquad \epsilon_n\ge0.
\]

\begin{proposition}[Remainders required of the training procedure]
\label{prop:algorithmic_remainders}
If $\sup_{\mathcal H_n}|R_n|=o_p(1)$ and $\epsilon_n=o_p(1)$, then
$\widetilde\eta_n$ is consistent.  Let
$\eta_n^*:=(\vartheta_0,q_n^*)$ be the sieve approximant in
Assumption~\ref{ass:sieve_estimation}.  If, additionally,
$R_n(\widetilde\eta_n)-R_n(\eta_n^*)=o_p(r_n^2)$ and
$\epsilon_n=o_p(r_n^2)$, then
$\rho(\widetilde\eta_n,\eta_0)=O_p(r_n)$.  If, additionally,
$\widetilde\vartheta_n$ is interior and lies in the profile neighborhood
$\|\widetilde\vartheta_n-\vartheta_0\|\le\delta_n$ with probability tending
to one, the perturbed nuisance/profile score equations have
$o_p(n^{-1/2})$ residuals, and
\[
 \sup_{\|\vartheta-\vartheta_0\|\le\delta_n}
 \|\nabla_\vartheta R_n^p(\vartheta)\|=o_p(n^{-1/2}),
 \qquad
 \sup_{\|\vartheta-\vartheta_0\|\le\delta_n}
 \|\nabla^2_{\vartheta\vartheta}R_n^p(\vartheta)\|=o_p(1),
\]
then Theorem~\ref{thm:sieve_clt} holds with
$\widetilde\vartheta_n$ in place of $\widehat\vartheta_n$.

For a simulated extremum criterion, suppose all preceding consistency and
profile-expansion conditions hold for the simulated criterion.  If $R$
conditionally i.i.d.\ draws per observation give score noise $\xi_r$ that is
conditionally mean zero given $U$ and has finite second moments, the fixed-$R$
covariance is
\[
 A_0^{-1}B_RA_0^{-1},\qquad
 B_R:=B_0+\frac1R\E_0[\mathrm{Var}(\xi_1\mid U)].
\]
Letting $R\to\infty$, or evaluating an analytic Gaussian ELBO, removes this
first-order simulation term under the corresponding triangular-array
conditions.  An objective gap of $o_p(r_n^2)$ alone does not imply the score
accuracy needed for the CLT; under smooth strong concavity, an
$o_p(n^{-1})$ objective gap is a typical sufficient order.  Generic
stochastic-gradient convergence to a global maximizer is not implied by the
statistical assumptions above.
\end{proposition}

\subsubsection{Objective and implementation scope}
\label{subsubsec:theory_code_scope}

The ELBO identity, and therefore every comparison theorem above, applies to
the implemented objective only in the genuine-VAE configuration: the
reconstruction and outcome terms are normalized log likelihoods, their
weights are one, the KL weight is one at the reported optimizer, free bits are
zero, annealing has terminated at one, and there is no unmodeled stochastic
decoder dropout.  WAE/MMD, deterministic AE, beta-VAE, active free bits,
powered prediction losses, and arbitrary MSE risks define different
population objectives and need their own estimating-equation and sandwich
theory.

Proposition~\ref{prop:mog_encoder_tightness} is an ideal-sieve existence
result, not a theorem about the stock finite network.  In the current
implementation, the component log variances are clamped below at $-8$, the
component and network widths do not grow with $n$, and a hard categorical
component draw omits the reconstruction/prediction gradient with respect to
mixture logits.  In addition, mixture-of-experts fusion is not a universal
joint multimodal encoder, and a fixed training budget and noisy checkpoint
criterion do not verify the global optimization remainder.  Instantiating the
theorem requires, at minimum, shrinkable component bandwidths, an explicit
$n$-indexed architecture, a valid mixture-component gradient (for example
Rao--Blackwellization), and optimization diagnostics at the required
per-record scale.

The current non-image encoder paths append prevalence covariates and optional
labels but omit content and prediction covariates.  By
Proposition~\ref{prop:omitted_conditioning_gap}, this generally prevents the
amortization gap from vanishing when those covariates affect the likelihood.
Under product fusion, appending the same optional labels to every modality
expert also repeats their evidence.  The repair is to omit outcomes from the
proposal or introduce one separate outcome expert exactly once.
The present Gaussian-prior switch also couples learning the prevalence mean
to learning a full covariance; an option that learns the mean with fixed
covariance is needed for the simplest normalized IdealPoint chart.

The stock subset routine draws one common active modality subset for a whole
minibatch and decodes only those active modalities.  It is therefore a
stochastic approximation to the expected-subset composite criterion in
Eq.~\eqref{eq:composite_subset_criterion}; the common draw correlates gradient
noise within the minibatch.  It is not the alternative likelihood based on one
independently realized coarsening per unit.  Per-unit subset draws are a
possible algorithmic refinement, not the behavior audited here.

Finally, model-matched simulations are essential.  The existing arithmetic
topic generator uses $p(w\mid\theta)=\theta^\top\Phi$, whereas the linear GTM
in Eq.~\eqref{eq:gtm_exact_linear_model} uses
$p(w\mid\theta)=\operatorname{softmax}(B\theta+Dx^c+b)$.  These are distinct
data-generating processes.  Likewise, posterior standard deviations of
unit-specific latent variables are not frequentist standard errors for global
generative parameters, and fixed within-unit information does not imply
pointwise consistency for realized latent positions.

Proofs are collected in Appendix~\ref{app:statistical_proofs}.
